% Apartado 8: limitaciones y siguientes pasos.
\needspace{8\baselineskip}
\section{Limitaciones y siguientes pasos}\label{sec:limitaciones}

El prototipo tiene límites que conozco. Los recojo aquí junto con lo que haría a continuación.

La primera son las cuotas gratuitas de Groq, de 8.000 tokens por minuto y una cuota diaria, las dos por modelo. Con
tráfico, eso se traduce en esperas de hasta 40~s por turno o en caer al modelo local, que es menos fiable. El
siguiente paso sería un plan de pago o un modelo propio en un servidor, y además la selección de herramientas por
categoría, para gastar menos tokens en cada ronda.

La segunda es el propio modelo local de 8B. Necesita guardas y respuestas fijas, y rinde peor en preguntas
complejas. Con una GPU mayor se podría usar un modelo local más grande, y las guardas ya son genéricas, así que
servirían igual.

La tercera es la licencia de los pesos de OmniVoice, CC-BY-NC, que no permite el uso comercial. Vale para un reto
académico, pero no para producción. Cambiarlo es sustituir el adaptador \cod{Sintetizador} por un modelo con licencia
comercial, que es una función en \ruta{composicion.py}.

La cuarta es la voz de entrada. El audio sale del equipo hacia Groq, que es un tercero. La solución sería ejecutar
Whisper en local, en la GPU.

La última es que no hay selección de herramientas. Cada ronda envía al modelo las 8, y con muchas más haría falta el
filtro por categoría en \cod{RegistroHerramientas} que describo en el apartado~\ref{sec:cuarenta}.
